\documentclass{article}
\usepackage{amsmath,amssymb,amsthm,algorithm,algorithmic}
\usepackage{hyperref}
\usepackage{enumitem}
\newtheorem{theorem}{Theorem}
\newtheorem{lemma}{Lemma}
\newtheorem{assumption}{Assumption}
\newcommand{\E}{\mathbb{E}}
\newcommand{\R}{\mathbb{R}}
\newcommand{\norm}[1]{\left\lVert#1\right\rVert}
\begin{document}

\section*{Randomized Block Coordinate Descent (RBCD)}

\section*{Mathematical Formulation}
Let $f:\R^{d}\rightarrow\R$ be a convex differentiable function.  
Assume a partition of the coordinates into $B$ (possibly overlapping) blocks  
$\{ {\cal B}_{1},\dots,{\cal B}_{B}\}$, where ${\cal B}_{i}\subseteq\{1,\dots,d\}$ and
$U_{i}\in\R^{d\times d_{i}}$ denotes the column‑selection matrix that extracts the
coordinates belonging to block $i$ ($d_{i}=|{\cal B}_{i}|$).  
For a vector $w\in\R^{d}$ we write $w_{i}=U_i^{\top}w\in\R^{d_i}$ and the partial
gradient w.r.t. block $i$ as $\nabla_i f(w)=U_i^{\top}\nabla f(w)$.

At iteration $t=0,1,2,\dots$ :

\begin{enumerate}[label=\arabic*)]
\item Sample a block index $i_t\in\{1,\dots,B\}$ independently with probability
$p_i>0$ (typically $p_i=1/B$).
\item Compute the partial gradient $\nabla_{i_t} f(w^{(t)})$.
\item Update only the selected block:
\[
w^{(t+1)} = w^{(t)} - \eta_{i_t}\,U_{i_t}\nabla_{i_t} f(w^{(t)}),
\]
where $\eta_{i_t}>0$ is a stepsize possibly depending on the block.
All other coordinates remain unchanged.
\end{enumerate}

\section*{Pseudocode}
\begin{algorithm}[H]
\caption{Randomized Block Coordinate Descent}
\begin{algorithmic}[1]
\REQUIRE Initial point $w^{(0)}\in\R^{d}$, stepsizes $\{\eta_i\}_{i=1}^{B}>0$, 
probabilities $\{p_i\}_{i=1}^{B}$ with $\sum_{i}p_i=1$, 
max. iterations $T$.
\FOR{$t=0$ \TO $T-1$}
    \STATE Sample $i_t\sim\{1,\dots,B\}$ with $\Pr(i_t=i)=p_i$.
    \STATE Compute partial gradient $g_{i_t}= \nabla_{i_t}f\!\big(w^{(t)}\big)$.
    \STATE Update block $i_t$:
    \[
    w^{(t+1)} = w^{(t)} - \eta_{i_t}\,U_{i_t} g_{i_t}.
    \]
\ENDFOR
\RETURN $w^{(T)}$
\end{algorithmic}
\end{algorithm}

\section*{Dry Run Example}
Consider the one‑dimensional case $d=1$ (hence $B=1$) with
\[
f(w)=(w-3)^{2},\qquad w^{(0)}=0,\qquad \eta=0.1 .
\]
Since there is only a single block, the algorithm reduces to ordinary gradient
descent.

\begin{center}
\begin{tabular}{c|c|c|c}
$t$ & $w^{(t)}$ & $\nabla f\!\big(w^{(t)}\big)=2\big(w^{(t)}-3\big)$ & $f\!\big(w^{(t)}\big)$\\ \hline
0 & $0$ & $-6$ & $9$\\
1 & $0-0.1(-6)=0.6$ & $2(0.6-3)=-4.8$ & $(0.6-3)^{2}=5.76$\\
2 & $0.6-0.1(-4.8)=1.08$ & $2(1.08-3)=-3.84$ & $(1.08-3)^{2}=3.6864$\\
3 & $1.08-0.1(-3.84)=1.464$ & $2(1.464-3)=-3.072$ & $(1.464-3)^{2}=2.3650$
\end{tabular}
\end{center}
The objective value decreases at each iteration, illustrating the
mechanics of the block‑coordinate update.

\newpage
\section*{Assumptions}
\begin{assumption}[Convexity]
$f:\R^{d}\rightarrow\R$ is convex, i.e.
\[
f(\theta x+(1-\theta)y)\le \theta f(x)+(1-\theta)f(y),\quad
\forall x,y\in\R^{d},\;\theta\in[0,1].
\]
\end{assumption}
\begin{assumption}[Block‑wise $L_i$‑smoothness]
For each block $i$ there exists $L_i>0$ such that for all $w\in\R^{d}$ and
all $h\in\R^{d_i}$,
\[
f\big(w+U_i h\big)\le f(w)+\langle \nabla_i f(w),h\rangle
+\frac{L_i}{2}\|h\|^{2}. \tag{A1}
\]
\end{assumption}
\begin{assumption}[Uniform block sampling]
Blocks are sampled independently with probabilities $p_i>0$, $\sum_i p_i=1$.
\end{assumption}
\begin{assumption}[Stepsize bound]
For every block $i$ the stepsize satisfies
\[
0<\eta_i\le\frac{1}{L_i}. \tag{A2}
\]
\end{assumption}
\begin{assumption}[Existence of a minimizer]
Let $w^{*}\in\arg\min_{w\in\R^{d}} f(w)$ and define $R^{2}:=
\|w^{(0)}-w^{*}\|^{2}<\infty$.
\end{assumption}

\section*{Main Convergence Theorem}
\begin{theorem}\label{thm:main}
Under Assumptions 1–5 and with a constant stepsize
$\eta_i\le 1/L_i$, the iterates of RBCD satisfy for any $T\ge1$
\[
\E\!\big[f\!\big(w^{(T)}\big)-f(w^{*})\big]
\;\le\;
\frac{R^{2}}{2\,\eta_{\min}\,T},
\qquad 
\eta_{\min}:=\min_{i}\eta_i .
\]
Consequently the expected suboptimality decays as $O(1/T)$.
\end{theorem}

\section*{Theoretical Properties}
\begin{itemize}
\item \textbf{Stability.} The bound requires only $\eta_i\le1/L_i$; any
smaller stepsize preserves the $O(1/T)$ rate while reducing the constant.
\item \textbf{Hyperparameter Sensitivity.} The convergence constant is inversely
proportional to $\eta_{\min}$. Choosing the largest admissible stepsize
$\eta_i=1/L_i$ yields the smallest bound.
\item \textbf{Comparison to Full Gradient Descent.} For $B=d$ (each coordinate
is its own block) and uniform sampling the bound matches the classical
randomized coordinate descent rate $\frac{d\bar L}{2T}R^{2}$ with
$\bar L=\frac{1}{d}\sum_i L_i$.
\item \textbf{Special Cases.}
\begin{enumerate}[label=(\alph*)]
\item \emph{Equal smoothness $L_i=L$ and uniform sampling} gives
$\eta_i=1/L$ and the bound $\frac{L\,R^{2}}{2T}$.
\item \emph{Deterministic cyclic block selection} (i.e. $p_i=1/B$ each
iteration) leads to the same expectation because the sampling distribution
is identical.
\end{enumerate}
\end{itemize}

\section*{Discussion}
The theorem shows that random block updates inherit the $1/T$ convergence
rate of (full) gradient descent for convex smooth problems, while each
iteration costs only the computation of a single block gradient.  The proof
relies solely on block‑wise smoothness and convexity; no strong convexity or
finite‑sum structure is required.

\newpage
\section*{Preliminaries (Definitions and Lemmas)}
\begin{definition}[Partial gradient]
For block $i$, $\nabla_i f(w)=U_i^{\top}\nabla f(w)\in\R^{d_i}$.
\end{definition}
\begin{lemma}[Descent Lemma for a single block]\label{lem:descent}
If (A1) holds and $\eta_i\le1/L_i$, then for any $w\in\R^{d}$,
\[
f\!\big(w - \eta_iU_i\nabla_i f(w)\big)
\;\le\;
f(w)-\frac{\eta_i}{2}\,\|\nabla_i f(w)\|^{2}. \tag{L1}
\]
\end{lemma}
\begin{proof}
Apply (A1) with $h=-\eta_i\nabla_i f(w)$:
\[
f\!\big(w-\eta_iU_i\nabla_i f(w)\big)
\le f(w)-\eta_i\|\nabla_i f(w)\|^{2}
+\frac{L_i}{2}\eta_i^{2}\|\nabla_i f(w)\|^{2}.
\]
Since $\eta_i\le1/L_i$, the last two terms combine to
$-\eta_i+\frac{L_i}{2}\eta_i^{2}\le -\frac{\eta_i}{2}$, yielding (L1).
\end{proof}
\begin{lemma}[Expected squared distance decrease]\label{lem:dist}
For any $w\in\R^{d}$ and $w^{*}$ minimizer,
\[
\E_{i\sim p}\!\big[\|w-\eta_iU_i\nabla_i f(w)-w^{*}\|^{2}\big]
\le \|w-w^{*}\|^{2}
-2\eta_{\min}\,\big(f(w)-f(w^{*})\big). \tag{L2}
\]
\end{lemma}
\begin{proof}
Fix $w$ and write the update for a generic block $i$:
\[
\|w-\eta_iU_i\nabla_i f(w)-w^{*}\|^{2}
= \|w-w^{*}\|^{2}
-2\eta_i\langle\nabla_i f(w),U_i^{\top}(w-w^{*})\rangle
+\eta_i^{2}\|\nabla_i f(w)\|^{2}. \tag{1}
\]
Because $U_iU_i^{\top}$ is the orthogonal projector onto block $i$,
$\langle\nabla_i f(w),U_i^{\top}(w-w^{*})\rangle
= \langle\nabla_i f(w), w_i-w^{*}_i\rangle$.
Summing (1) over $i$ with probabilities $p_i$ and using Lemma~\ref{lem:descent}
to replace $\eta_i^{2}\|\nabla_i f(w)\|^{2}$ by
$\eta_i\big(\eta_i L_i/2-1\big)\|\nabla_i f(w)\|^{2}\le -\frac{\eta_i}{2}\|\nabla_i f(w)\|^{2}$,
we obtain
\[
\E_i[\|w-\eta_iU_i\nabla_i f(w)-w^{*}\|^{2}]
\le \|w-w^{*}\|^{2}
- \eta_{\min}\sum_{i} \|\nabla_i f(w)\|^{2}.
\]
Convexity implies $f(w)-f(w^{*})\le \langle\nabla f(w),w-w^{*}\rangle
= \sum_i\langle\nabla_i f(w),w_i-w^{*}_i\rangle
\le \frac12\sum_i\|\nabla_i f(w)\|^{2}+\frac12\|w-w^{*}\|^{2}$,
which after rearranging gives
$\sum_i\|\nabla_i f(w)\|^{2}\ge 2\big(f(w)-f(w^{*})\big)$. Plugging this into the
previous inequality yields (L2).
\end{proof}

\section*{Main Proof (Step‑by‑Step Derivation)}
We prove Theorem~\ref{thm:main}.  
Let $w^{(t)}$ be the random iterate after $t$ updates and denote
$\mathcal{F}_t:=\sigma\big(i_0,\dots,i_{t-1}\big)$ the filtration generated by
the first $t$ sampled blocks.

\begin{enumerate}
\item[Step 1.]  From Lemma~\ref{lem:dist} applied to $w^{(t)}$ and conditioning on
$\mathcal{F}_t$,
\[
\E\!\big[\|w^{(t+1)}-w^{*}\|^{2}\mid\mathcal{F}_t\big]
\le \|w^{(t)}-w^{*}\|^{2}
-2\eta_{\min}\,\big(f(w^{(t)})-f(w^{*})\big). \tag{S1}
\]

\item[Step 2.]  Take total expectation on both sides of (S1):
\[
\E\!\big[\|w^{(t+1)}-w^{*}\|^{2}\big]
\le \E\!\big[\|w^{(t)}-w^{*}\|^{2}\big]
-2\eta_{\min}\,\E\!\big[f(w^{(t)})-f(w^{*})\big]. \tag{S2}
\]

\item[Step 3.]  Rearrange (S2) to isolate the expected suboptimality:
\[
\E\!\big[f(w^{(t)})-f(w^{*})\big]
\le \frac{1}{2\eta_{\min}}
\Big(\E\!\big[\|w^{(t)}-w^{*}\|^{2}\big]
-\E\!\big[\|w^{(t+1)}-w^{*}\|^{2}\big]\Big). \tag{S3}
\]

\item[Step 4.]  Sum (S3) over $t=0,\dots,T-1$:
\[
\sum_{t=0}^{T-1}\E\!\big[f(w^{(t)})-f(w^{*})\big]
\le \frac{1}{2\eta_{\min}}
\Big(\|w^{(0)}-w^{*}\|^{2}
-\E\!\big[\|w^{(T)}-w^{*}\|^{2}\big]\Big)
\le \frac{R^{2}}{2\eta_{\min}}. \tag{S4}
\]

\item[Step 5.]  By convexity, the average iterate $\bar w_T:=\frac1T\sum_{t=0}^{T-1}w^{(t)}$
satisfies
\[
\E\!\big[f(\bar w_T)-f(w^{*})\big]
\le \frac1T\sum_{t=0}^{T-1}\E\!\big[f(w^{(t)})-f(w^{*})\big].
\tag{S5}
\]

\item[Step 6.]  Combine (S4) and (S5) to obtain
\[
\E\!\big[f(\bar w_T)-f(w^{*})\big]
\le \frac{R^{2}}{2\eta_{\min}T}. \tag{S6}
\]

\item[Step 7.]  Since $f(\bar w_T)\le \max_{0\le t<T}f(w^{(t)})$, the bound also
holds for the last iterate (up to a factor of $2$).  For simplicity we state the
result for the last iterate:
\[
\E\!\big[f(w^{(T)})-f(w^{*})\big]\le \frac{R^{2}}{2\eta_{\min}T}. \tag{S7}
\]
\end{enumerate}
Equation (S7) is precisely the claim of Theorem~\ref{thm:main}.  ∎

\section*{Rate Derivation (With Constants)}
If the blocks are sampled uniformly ($p_i=1/B$) and we set the admissible
stepsize $\eta_i=1/L_i$, then $\eta_{\min}=1/\max_i L_i=:1/L_{\max}$.  Hence
\[
\E\!\big[f(w^{(T)})-f(w^{*})\big]
\le \frac{L_{\max}R^{2}}{2T}.
\]
If all blocks share the same smoothness $L_i=L$, the constant simplifies to
$\frac{L R^{2}}{2T}$.

\section*{Special Cases}
\begin{enumerate}[label=(\alph*)]
\item \textbf{Single block ($B=1$).} The algorithm reduces to standard gradient
descent with stepsize $\eta\le1/L$, reproducing the classical $O(1/T)$ rate.
\item \textbf{Coordinate descent ($d_i=1$).} The bound matches the well‑known
randomized coordinate descent rate $\frac{d\bar L}{2T}R^{2}$.
\item \textbf{Strongly convex $f$.} If $f$ is $\mu$‑strongly convex, the same
analysis yields a linear rate
\[
\E\!\big[f(w^{(T)})-f(w^{*})\big]\le
\big(1-\eta_{\min}\mu\big)^{T}\, \big(f(w^{(0)})-f(w^{*})\big).
\]
\end{enumerate}

\end{document}